% Chapter 20: Trees and finite differences.

Monte Carlo samples paths forward in time. Grids compute the same expectation backwards:
they discretize the state, replace the conditional expectation over one time step by a
weighted average of neighbouring values, and step from the payoff to the price today. A
\emph{tree} takes the weights from a discrete-time model that converges to the diffusion; a
\emph{finite-difference scheme} from difference quotients in the pricing equation. The two meet
in the explicit scheme, a trinomial tree up to its discounting. Grids give the price for all spots at once,
the Greeks almost for free, and early exercise for one maximum per node. This is why \cref{ch:14-american-mc} took its American reference values
from a smoothed binomial tree, studied in \cref{sec:fd-binomial}, and its Bermudan ones from a
backward induction on a grid in log price with the exact Gaussian transition density, a
quadrature cousin of the schemes of \cref{sec:fd-schemes}. Their cost grows exponentially with the number
of state variables (\cref{prop:mc-dimension}), so they are the method of choice in one or two
dimensions and of no use in ten.

Throughout, $S$ is the stock with yield $y$ of \cref{def:sm-model} under $\Q$, with log-drift
$\nu=r-y-\sigma^2/2$, and $r\ge0$. The price at time to maturity $\tau$ of the European payoff
$g(S_T)$, as a function of the log price $x=\log S$, is
$u(\tau,x)=e^{-r\tau}\EQ\big[g\big(e^{x+\nu\tau+\sigma\sqrt\tau Z}\big)\big]$ with $Z\sim\mathcal N(0,1)$, and by
\cref{thm:sm-feynman-kac} (whose code checks it with difference quotients) it solves
\begin{equation}\label{eq:fd-pde}
  \partial_\tau u=\tfrac12\sigma^2\partial_{xx}u+\nu\,\partial_xu-ru,\qquad u(0,x)=g(e^x),
\end{equation}
whenever $g$ is regular enough for the equation to make sense. The running example is the
course's test case $S_0=K=100$, $T=1$, $r=1\%$, $y=2\%$, $\sigma=20\%$, with call $7.36429$
and put $8.34941$.

\section{Binomial trees}\label{sec:fd-binomial}

\begin{definition}[Binomial tree]\label{def:fd-tree}
Fix $n\ge1$ steps of length $\Delta t=T/n$ and factors $0<d<u$. The tree has the nodes
$S_0u^jd^{i-j}$, $0\le j\le i\le n$, at time $i\Delta t$, and the \emph{risk-neutral probability}
$q=(e^{(r-y)\Delta t}-d)/(u-d)$. The values $V_n(j)=g(S_0u^jd^{n-j})$ are propagated backwards by
\begin{equation}\label{eq:fd-tree}
  V_i(j)=e^{-r\Delta t}\big[q\,V_{i+1}(j+1)+(1-q)\,V_{i+1}(j)\big],
\end{equation}
and the tree price is $V_0(0)$; for an American option each $V_i(j)$ is replaced by its maximum
with the payoff $g(S_0u^jd^{i-j})$. The \emph{Cox--Ross--Rubinstein} (CRR) tree has $u=e^{\sigma\sqrt{\Delta t}}$ and
$d=1/u$ \citep{crr1979}; the Jarrow--Rudd tree has $u,d=e^{\nu\Delta t\pm\sigma\sqrt{\Delta t}}$ \citep{jarrow-rudd-1983}.
\end{definition}

The tree is a model in its own right when $q\in(0,1)$, that is when $d<e^{(r-y)\Delta t}<u$: this is
the no-arbitrage condition of the one-period model, and $V_i(j)$ is then the price in the
binomial model whose stock moves by the factor $u$ or $d$ and pays dividends at the rate $y$
(the IEOR 222 notes treat $y=0$). For the European payoff, \eqref{eq:fd-tree} unrolls into
$V_0(0)=e^{-rT}\sum_{j=0}^n\binom nj q^j(1-q)^{n-j}g(S_0u^jd^{n-j})$, which our code evaluates in $O(n)$
operations. For the CRR tree the terminal log prices are $\log S_0+(2j-n)\delta$ with
$\delta=\sigma\sqrt{T/n}$, a lattice of spacing $2\delta$, and we write
\[
  \theta_n=\operatorname{frac}\Big(\frac{\log(K/S_0)+n\delta}{2\delta}\Big)\in[0,1)
\]
for the position of $\log K$ between two consecutive terminal nodes: $\theta_n=0$ when $K$ is a
node. At $S_0=K$, $\theta_n=0$ for even $n$ and $\theta_n=\frac12$ for odd $n$.

\begin{proposition}[Binomial trees]\label{prop:fd-binomial}
Let $\sigma>0$ and let $C_n$ be the CRR price of a European payoff $g$.
\begin{enumerate}[label=(\roman*)]
  \item \emph{Validity.} $q\in(0,1)$ if and only if $n>(r-y)^2T/\sigma^2$; if $r=y$, $q=1/(1+u)$ for
    every $n$. As $n\to\infty$, $q=\frac12+\frac{\nu}{2\sigma}\sqrt{\Delta t}+O(\Delta t)$.
  \item \emph{Convergence.} If $g$ is continuous with $|g(s)|\le a(1+s)$, then
    $C_n\to e^{-rT}\EQ[g(S_T)]$. The same holds for a bounded $g$ whose discontinuities form a
    finite set (digitals).
  \item \emph{Oscillation.} For the call with strike $K>0$, $C_n-C=O(1/n)$, where $C$ is the
    Black--Scholes price, but $n(C_n-C)$ has no limit in general:
    \begin{equation}\label{eq:fd-oscillation}
      n\,(C_n-C)=A+B\,\theta_n(1-\theta_n)+o(1),\qquad B=2\sigma\sqrt T\,Ke^{-rT}N'(d_-),
    \end{equation}
    with a constant $A$ depending on $(S_0,K,T,r,y,\sigma)$ only.
  \item \emph{Better trees.} (a) The BBS tree of \textcite{broadie-detemple} replaces the values one
    step before maturity by the Black--Scholes prices with maturity $\Delta t$; BBSR is the
    extrapolation $2C^{\mathrm{BBS}}_n-C^{\mathrm{BBS}}_{n/2}$. (b) The tree of
    \textcite{leisen-reimer-1996} takes $n$ odd, $q=h_n^{-1}(d_-)$, $q'=h_n^{-1}(d_+)$ with the
    Peizer--Pratt inversion $h_n^{-1}(z)=\frac12+\frac12\operatorname{sign}(z)
    \big[1-\exp\big(-\big(\frac{z}{n+1/3+0.1/(n+1)}\big)^2\big(n+\frac16\big)\big)\big]^{1/2}$
    \citep{peizer-pratt-1968}, $u=e^{(r-y)\Delta t}q'/q$ and $d=(e^{(r-y)\Delta t}-qu)/(1-q)$. Its call
    prices converge at order $1/n^2$, without oscillation.
\end{enumerate}
\end{proposition}

\begin{proof}
(i) $q\in(0,1)$ if and only if $e^{-\sigma\sqrt{\Delta t}}<e^{(r-y)\Delta t}<e^{\sigma\sqrt{\Delta t}}$, that is
$|r-y|\Delta t<\sigma\sqrt{\Delta t}$, that is $n>(r-y)^2T/\sigma^2$. If $r=y$,
$q=(1-u^{-1})/(u-u^{-1})=1/(1+u)$. With $\delta=\sigma\sqrt{\Delta t}$, the numerator of $q$ is
$\delta-\frac12\delta^2+(r-y)\Delta t+O(\Delta t^{3/2})=\delta+\nu\Delta t+O(\Delta t^{3/2})$ and the denominator is
$2\delta+O(\delta^3)$.

(ii) $\log(S_n/S_0)=X_n=\xi_1+\dots+\xi_n$ with independent $\xi_i=\pm\delta$, $\Q(\xi_i=\delta)=q$. By (i),
$\E[\xi_i]=(2q-1)\delta=\nu\Delta t+O(\Delta t^{3/2})$ and $\E[\xi_i^2]=\sigma^2\Delta t$, so, as $|\xi_i|\le\delta$,
$\E[e^{is\xi_i}]=1+is\nu\Delta t-\frac12s^2\sigma^2\Delta t+O(\Delta t^{3/2})$ for each real $s$, and
$\E[e^{isX_n}]=\E[e^{is\xi_1}]^n\to e^{is\nu T-s^2\sigma^2T/2}$. By Lévy's continuity theorem $X_n$ converges in
law to $\mathcal N(\nu T,\sigma^2T)$, so $S_n$ to $S_T$. Likewise
$\E[S_n^2]=S_0^2\big(qu^2+(1-q)u^{-2}\big)^n=S_0^2\big(1+(2\nu+2\sigma^2)\Delta t+O(\Delta t^{3/2})\big)^n$ is
bounded in $n$, so $(g(S_n))_n$ is bounded in $L^2$, hence uniformly integrable. By the mapping
theorem, $g(S_n)$ converges in law to $g(S_T)$ when the discontinuities of $g$ have probability zero
under the law of $S_T$, which has a density \citep[Theorem~2.7]{billingsley}; uniform
integrability then gives convergence of the means \citep[Theorem~3.5]{billingsley}.

(iii) and (iv) are cited. \textcite{diener-diener-2004} compute the expansion for the CRR tree;
\textcite{walsh-2003}, Corollary~4.1, proves \eqref{eq:fd-oscillation} with this $B$ (his
$2\sigma^2T\tilde K\hat p$, with $\tilde K=Ke^{-rT}$) for $n$ even in the CRR tree of the discounted price, where $\theta_n$
is the position of the discounted strike. For (iv)(b),
\textcite{leisen-reimer-1996} observed the order $1/n^2$ and conjectured it, and
\textcite{joshi-2010} proved it; part (iv)(a) is a definition, discussed after
\cref{fig:fd-trees}.
\end{proof}

\begin{remark}[Where the oscillation comes from]\label{rem:fd-oscillation}
By a local central limit theorem the tree puts mass close to $H\,f(x_j)$ on the terminal node $x_j$,
where $H=2\delta$ and $f$ is a smooth density close to that of $\log S_T$. The call is then
$e^{-rT}H\sum_jF(x_j)$ with $F(x)=f(x)(e^x-K)^+$, a function that vanishes at $\log K$ with slope
$Kf(\log K)$ to the right. The Euler--Maclaurin formula for a lattice whose first point lies $\alpha H$
above $c$, $H\sum_{m\ge0}F(c+(m+\alpha)H)=\int_c^\infty F-HB_1(\alpha)F(c)-\frac{H^2}2B_2(\alpha)F'(c)+O(H^3)$, with the
Bernoulli polynomials $B_1(\alpha)=\alpha-\frac12$ and $B_2(\alpha)=\alpha^2-\alpha+\frac16$, applied at $c=\log K$ with
$\alpha=1-\theta_n$ and $B_2(1-\theta_n)=B_2(\theta_n)$, leaves $e^{-rT}\frac{H^2}2\theta_n(1-\theta_n)Kf(\log K)$ besides terms that do not depend
on $\theta_n$; with $H^2=4\sigma^2T/n$ and $f(\log K)=N'(d_-)/(\sigma\sqrt T)$ this is the $B$ of
\eqref{eq:fd-oscillation}. The argument does not control the error of the local expansion
(Walsh's proof does, by Skorokhod embedding), but it shows where $B$ comes from, and it applies to
the CRR tree as defined here: on the test case $B=15.6222$, and fitting $n(C_n-C)$ on $1$ and
$\theta_n(1-\theta_n)$ over $n=1000,\dots,2000$ gives $A=-1.964$, $B=15.6222$; for $K=110$, $B=14.2815$ against
fits of $14.2859$ ($n=1000,\dots,2000$) and $14.2800$ ($n=8000,\dots,8400$). For a digital the same computation leaves the
first-order term $-HB_1(\alpha)F(c)$: its error oscillates at order $n^{-1/2}$ (\exref{ex:fd-digital}).
\end{remark}

\begin{figure}[ht]
  \centering
  \includegraphics[width=\textwidth]{ch20-trees}
  \caption{CRR trees on the test case. (a) $n(C_n-C)$ for the calls with $K=100$ (two branches:
    $n$ even and odd) and $K=110$; the black marks are $A+B\theta_n(1-\theta_n)$ with $B$ from
    \eqref{eq:fd-oscillation} and $A$ fitted on $n=1000,\dots,2000$. (b) Absolute errors of the European
    call: CRR for $n$ even and odd, BBS for $n$ even, BBSR and Leisen--Reimer, and the trinomial
    tree of \cref{sec:fd-trinomial} with $\kappa=\sqrt{3/2}$; $n$ runs over multiples of $4$
    between $20$ and $4000$ ($n+1$ for the odd-$n$ trees; up to $2000$ for the trinomial tree).}
  \label{fig:fd-trees}
\end{figure}

\paragraph{Parity, smoothing, extrapolation.} At $S_0=K$ the two parities give two smooth
sequences, $nC_n-nC\to A=-1.964$ for $n$ even and $A+B/4=1.942$ for $n$ odd: the CRR price is below the
Black--Scholes price for even $n$ and above it for odd $n$ (\cref{fig:fd-trees}(a)), and
Richardson extrapolation across consecutive $n$ fails. Away from the money $\theta_n$ wanders
through $[0,1)$, and doubling $n$ can increase the error. The remedies of (iv) attack the kink of
the payoff. The BBS tree, whose last step is smooth on the scale of the lattice, does not remove
the oscillation, but on the test case it shrinks the parity effect
from $B/4=3.906$ to $0.023$ in units of $1/n$ ($nC^{\mathrm{BBS}}_n-nC$ is $0.6498$ at $n=400$ and
$0.6266$ at $n=401$). At the money BBSR is then of order $1/n^2$ for $n$ divisible by $4$ (so that $n$
and $n/2$ have the same parity) and of order $1/n$ for $n\equiv2\pmod4$ ($n$ times its error is $0.061$
at $n=102$, $0.047$ at $n=1602$). For $K=110$, $n$ times the BBS error moves between $0.467$ and $0.488$,
and BBSR is first order ($n^2$ times its error: $5.8$, $-1.6$, $-28.0$ at $n=400$, $800$, $1600$). The Leisen--Reimer tree places the strike between
the two middle terminal nodes for every odd $n$ (we checked this for strikes from $80$ to $130$)
and chooses $q$ and $q'$ so that the binomial probabilities of finishing above the strike, under
the two measures of the call formula, approximate $N(d_-)$ and $N(d_+)$: $n^2(C_n-C)=-0.3955$ at
$n=201$ and $-0.3978$ at $n=801$ (\cref{fig:fd-trees}(b)). Neither survives early exercise
unchanged: on the first Longstaff--Schwartz put of \cref{tab:fd-american}, the Leisen--Reimer error
halves when $n$ doubles ($-4.8\times10^{-3}$ at $n=101$, $-2.9\times10^{-4}$ at $n=1601$), because the American
price is not smooth across the exercise boundary.

\paragraph{American options on trees.} With the projection $V_i(j)\leftarrow\max(V_i(j),g)$ the
CRR tree converges to the American price; for the put, \textcite{lamberton-2018}, Theorem~1.1 (arXiv version),
bounds the error by a constant times $(\log n)^{3/2}/n$ ($(\log n)/n$ if $y>r$). The American
references of \cref{ch:14-american-mc} are BBSR trees (its code and ours agree to rounding);
\cref{tab:fd-american} compares them with the grids of \cref{sec:fd-american}.

\section{Trinomial trees and the explicit scheme}\label{sec:fd-trinomial}

A trinomial tree moves the log price by $+h$, $0$ or $-h$ in each step. Following
\textcite{kamrad-ritchken-1991} (and the three-jump process of \textcite{boyle-1986}), take
$h=\kappa\sigma\sqrt{\Delta t}$ with a \emph{stretch} $\kappa>0$ and
\begin{equation}\label{eq:fd-kr}
  p_\pm=\frac1{2\kappa^2}\pm\frac{\nu\sqrt{\Delta t}}{2\kappa\sigma},\qquad p_0=1-\frac1{\kappa^2},
\end{equation}
with backward induction $V\leftarrow e^{-r\Delta t}(p_+V_++p_0V_0+p_-V_-)$ on the nodes
$\log S_0+mh$, $|m|\le i$, at step $i$. The tree needs no boundary condition: in $n$ steps it reaches
$n$ nodes away.

\begin{proposition}[Trinomial trees]\label{prop:fd-trinomial}
Let $\kappa>0$, $\Delta t=T/n$ and $h=\kappa\sigma\sqrt{\Delta t}$.
\begin{enumerate}[label=(\roman*)]
  \item \emph{Moments and positivity.} An increment $\xi$ has $\E[\xi]=\nu\Delta t$ and
    $\E[\xi^2]=\sigma^2\Delta t$ exactly. The three weights are probabilities if and only if $\kappa\ge1$ and
    $h|\nu|\le\sigma^2$. If $\nu=0$, then $\E[\xi^4]=\kappa^2\sigma^4\Delta t^2$, which equals the Gaussian value
    $3\sigma^4\Delta t^2$ if and only if $\kappa=\sqrt3$, and then $(p_+,p_0,p_-)=(\frac16,\frac23,\frac16)$.
  \item \emph{The explicit scheme.} The explicit finite-difference scheme for \eqref{eq:fd-pde}
    with time step $k=\Delta t$ and mesh $h$ (\cref{sec:fd-schemes}) is
    $U^{n+1}_j=p_+U^n_{j+1}+p_0U^n_j+p_-U^n_{j-1}-rk\,U^n_j$ with the weights \eqref{eq:fd-kr}: the tree with
    the discount $e^{-rk}$ replaced by a subtraction of $rkU^n_j$, which differs by $O(k^2)$ per step.
    Its parameter $\sigma^2k/(2h^2)$ equals $1/(2\kappa^2)$, so $\kappa\ge1$ is the stability condition
    $\sigma^2k\le h^2$ of \cref{thm:fd-schemes}(i).
  \item \emph{Convergence.} If $\kappa\ge1$ and $h|\nu|\le\sigma^2$, the tree is a positive operator
    with $\|V\|_\infty$ nonincreasing from step to step, and if $\partial_\tau^iu$ ($i\le2$) and
    $\partial_x^ju$ ($j\le4$) are continuous and bounded on $[0,T]\times\R$, its error is at most
    $c\,T\Delta t$ uniformly in the nodes, with $c$ depending on these bounds and on $\kappa$. If $\kappa<1$, the
    backward induction multiplies the alternating sequence $(-1)^m$ by a factor tending to
    $2/\kappa^2-1>1$ in modulus as $\Delta t\to0$.
\end{enumerate}
\end{proposition}

\begin{proof}
(i) $\E[\xi]=h(p_+-p_-)=\kappa\sigma\sqrt{\Delta t}\,\nu\sqrt{\Delta t}/(\kappa\sigma)$ and $\E[\xi^2]=h^2(p_++p_-)=h^2/\kappa^2$.
$p_0\ge0$ if and only if $\kappa\ge1$, and $p_\pm\ge0$ if and only if $\sigma\ge\kappa|\nu|\sqrt{\Delta t}$, that is
$\sigma^2\ge|\nu|h$. If $\nu=0$, $\E[\xi^4]=h^4/\kappa^2=\kappa^2\sigma^4\Delta t^2$.
(ii) With $p_\pm^{\mathrm{fd}}=\sigma^2/(2h^2)\pm\nu/(2h)$ and $p_0^{\mathrm{fd}}=-\sigma^2/h^2-r$ (\cref{sec:fd-schemes}), the
explicit step is $U^{n+1}=U^n+k(p_-^{\mathrm{fd}}U^n_{j-1}+p_0^{\mathrm{fd}}U^n_j+p_+^{\mathrm{fd}}U^n_{j+1})$, and
$kp^{\mathrm{fd}}_\pm=\sigma^2\Delta t/(2\kappa^2\sigma^2\Delta t)\pm\nu\Delta t/(2\kappa\sigma\sqrt{\Delta t})=p_\pm$, $1+kp^{\mathrm{fd}}_0=p_0-rk$.
(iii) Under the conditions the weights are nonnegative with sum $e^{-r\Delta t}\le1$, which gives the
first claims. Taylor's formula with $\E[\xi]$, $\E[\xi^2]$ from (i), $|\E[\xi^3]|=h^2|\E[\xi]|$ and
$\E[\xi^4]=h^2\sigma^2\Delta t$ gives $e^{-r\Delta t}\E[u(\tau,x+\xi)]=u(\tau+\Delta t,x)+O(\Delta t^2)$, with a constant
depending on the bounds; a positive step with weights summing to at most one adds these local
errors without amplifying them, and $n\Delta t=T$. If $\kappa<1$, one step maps $(-1)^m$ to
$\big(e^{-r\Delta t}(p_0-p_+-p_-)\big)(-1)^m$, and $p_0-p_+-p_-=1-2/\kappa^2<-1$.
\end{proof}

The identity (ii) is the synthesis of \textcite{brennan-schwartz-1978}: an explicit scheme is a
jump process (our tests check it to $10^{-11}$). For $\kappa=0.9$ the backward induction returns
$-3\times10^{64}$ for the call after $400$ steps. On the test case $n(C_n-C)$ stays within $0.01$ of
$-1.97$, $-0.53$ and $-2.06$ for $\kappa=1$, $\sqrt{3/2}$ and $\sqrt3$ and $n$ from $400$ to $1600$ (the strike
stays a node, so there is no oscillation): the choice $\kappa=\sqrt3$ of the course's trinomial walk, best for one step of a smooth
payoff (see the exercises of \cref{ch:01-standard-model}), is not the best for a kink.

\section{Finite differences: stability and order}\label{sec:fd-schemes}

A finite-difference grid separates what a tree ties together: the mesh $h$ and the time step
$k$. Take $J$ intervals of a domain $[x_0,x_J]$ in log price, nodes $x_j=x_0+jh$, times $\tau_n=nk$,
$N=T/k$, and replace the derivatives in \eqref{eq:fd-pde} by central differences:
\[
  (L_hV)_j=p_-V_{j-1}+p_0V_j+p_+V_{j+1},\qquad p_\pm=\frac{\sigma^2}{2h^2}\pm\frac{\nu}{2h},\qquad p_0=-\frac{\sigma^2}{h^2}-r .
\]

\begin{definition}[The $\theta$-scheme]\label{def:fd-theta}
For $\theta\in[0,1]$, the $\theta$-scheme for \eqref{eq:fd-pde} is
\begin{equation}\label{eq:fd-theta}
  \frac{U^{n+1}_j-U^n_j}{k}=\theta\,(L_hU^{n+1})_j+(1-\theta)\,(L_hU^n)_j,\qquad 0<j<J,
\end{equation}
with $U^0_j=g(e^{x_j})$ and boundary values $U^n_0$, $U^n_J$ (\cref{sec:fd-boundary}). It is
\emph{explicit} for $\theta=0$, \emph{implicit} for $\theta=1$, and \emph{Crank--Nicolson} for $\theta=\frac12$
\citep{crank-nicolson-1947}. We write $\lambda=\sigma^2k/(2h^2)$. In matrix form,
$BV^{n+1}=CV^n+kF^n$ with $B=I-\theta kA$, $C=I+(1-\theta)kA$, $A$ tridiagonal with rows $(p_-,p_0,p_+)$, and
$F^n$ carrying the boundary values; the \emph{amplification operator} is $G=B^{-1}C$.
\end{definition}

For $\theta>0$ each step solves a tridiagonal system in $O(J)$ operations. To separate the
discretization from the truncation of the domain, the theorem considers the scheme on the
infinite grid $h\Z$, with the norms $\|V\|_\infty=\sup_j|V_j|$ and $\|V\|_{2,h}=(h\sum_jV_j^2)^{1/2}$. The scheme
is \emph{stable} in a norm, along a family of grids with $k\to0$, if there is $C_s$ with
$\|G^m\|\le C_s$ and $\|G^mB^{-1}\|\le C_s$ for all $mk\le T$; the \emph{truncation error} is
$T^n_j=(u^{n+1}_j-u^n_j)/k-\theta(L_hu^{n+1})_j-(1-\theta)(L_hu^n)_j$ with $u^n_j=u(\tau_n,x_j)$.

\begin{theorem}[Explicit, implicit and Crank--Nicolson schemes]\label{thm:fd-schemes}
Let $\sigma>0$, $r\ge0$ and $\theta\in[0,1]$, and consider the $\theta$-scheme on $h\Z$.
\begin{enumerate}[label=(\roman*)]
  \item \emph{Von Neumann stability.} The scheme maps $e^{\mathrm i j\xi h}$ to $g(\xi)e^{\mathrm i j\xi h}$, with
    \begin{equation}\label{eq:fd-symbol}
      g(\xi)=\frac{1+(1-\theta)z}{1-\theta z},\qquad
      z=k\Big(-\frac{2\sigma^2}{h^2}\sin^2\frac{\xi h}2-r+\mathrm i\frac{\nu}{h}\sin\xi h\Big),
    \end{equation}
    and it is stable in $\|\cdot\|_{2,h}$ if and only if $|g(\xi)|\le1+c\,k$ for all $\xi$ and some $c$. If
    $\theta\ge\frac12$, $|g(\xi)|\le1$ for all $\xi$, $k$ and $h$: implicit and Crank--Nicolson are
    \emph{unconditionally} stable in $\|\cdot\|_{2,h}$. If $\theta<\frac12$, along grids with $\lambda$ fixed the
    scheme is stable if and only if $(1-2\theta)\lambda\le\frac12$; the explicit scheme if and only if
    $\sigma^2k\le h^2$.
  \item \emph{Consistency.} If $\partial^i_\tau u$ ($i\le3$) and $\partial^j_xu$ ($j\le4$) are continuous and bounded by $M$
    on $[0,T]\times\R$, then
    $|T^n_j|\le M\big[(\frac12+\theta)k+(\frac{\sigma^2}{24}+\frac{|\nu|}6)h^2\big]$, and for $\theta=\frac12$,
    $|T^n_j|\le M\big[\frac16k^2+(\frac{\sigma^2}{24}+\frac{|\nu|}6)h^2\big]$: orders $O(k+h^2)$ and $O(k^2+h^2)$.
  \item \emph{Convergence.} If the scheme is stable in a norm, then
    $\|u^n-U^n\|\le C_sT\max_{m<n}\|T^m\|$ for $nk\le T$. Conversely, a consistent scheme for a well-posed
    linear problem that is not stable does not converge for every initial value
    \citep{lax-richtmyer-1956}.
  \item \emph{Positivity.} If $h|\nu|\le\sigma^2$ and $(1-\theta)k(\sigma^2/h^2+r)\le1$, then $B^{-1}$ and $G$ map
    nonnegative sequences to nonnegative sequences, $\|B^{-1}\|_\infty\le1$ and $\|G\|_\infty\le1$. Without
    the second condition positivity can fail: for Crank--Nicolson applied to
    $\partial_\tau u=\frac12\sigma^2\partial_{xx}u$, $G$ is positive if and only if $\lambda\le\frac32$.
\end{enumerate}
\end{theorem}

\begin{proof}
(i) $(L_he^{\mathrm i\cdot\xi h})_j=\big(\frac{\sigma^2}{2h^2}(e^{\mathrm i\xi h}-2+e^{-\mathrm i\xi h})+\frac{\nu}{2h}(e^{\mathrm i\xi h}-e^{-\mathrm i\xi h})-r\big)e^{\mathrm ij\xi h}$,
which is $(z/k)e^{\mathrm ij\xi h}$; solve \eqref{eq:fd-theta} for $U^{n+1}$. For $V\in\ell^2(\Z)$, the Fourier series
$\hat V(\xi)=h\sum_jV_je^{-\mathrm ij\xi h}$ satisfies Parseval's identity $\|V\|_{2,h}^2=\frac1{2\pi}\int_{-\pi/h}^{\pi/h}|\hat V|^2$, and
the scheme multiplies $\hat V$ by $g$, so $\|G^m\|_{2,h}=\max_\xi|g(\xi)|^m$. If $|g|\le1+ck$ this is at most
$e^{cT}$ for $mk\le T$; conversely, $\max|g|^{\lfloor T/k\rfloor}\le C_s$ forces $\max|g|\le C_s^{1/\lfloor T/k\rfloor}=1+O(k)$. Also
$B^{-1}$ multiplies by $1/(1-\theta z)$, of modulus at most $1$ because $\operatorname{Re}z\le0$. Write $z=x+\mathrm iw$. Then
\[
  |1+(1-\theta)z|^2-|1-\theta z|^2=2x+(1-2\theta)|z|^2 .
\]
If $\theta\ge\frac12$, this is at most $2x\le0$, so $|g|\le1$. Let $\theta<\frac12$ and $s=\sin^2(\xi h/2)$, so that
$x=-4\lambda s-rk$ and $w^2=(\nu k/h)^2\,4s(1-s)\le2\nu^2\lambda k/\sigma^2$. If $(1-2\theta)\lambda\le\frac12$, then
$(1-2\theta)(4\lambda s+rk)\le2s+rk$ and
\[
  2x+(1-2\theta)|z|^2\le(4\lambda s+rk)(2s+rk-2)+w^2\le(4\lambda+rk)rk+2\nu^2\lambda k/\sigma^2,
\]
while $|1-\theta z|\ge1$; so $|g|^2\le1+c'k$. If $(1-2\theta)\lambda>\frac12$, at $\xi h=\pi$ we have $w=0$ and
$g=\frac{1-(1-\theta)(4\lambda+rk)}{1+\theta(4\lambda+rk)}\to\frac{1-4(1-\theta)\lambda}{1+4\theta\lambda}<-1$ as $k\to0$, so $|g|\le1+ck$ fails.

(ii) By Taylor's formula, $\frac1{h^2}(f(x+h)-2f(x)+f(x-h))=f''(x)+\frac{h^2}{12}f''''(\eta)$ and
$\frac1{2h}(f(x+h)-f(x-h))=f'(x)+\frac{h^2}6f'''(\eta')$ for some $\eta,\eta'$, so $L_hu$ differs from
$\frac12\sigma^2\partial_{xx}u+\nu\partial_xu-ru=\partial_\tau u$ by at most $Mh^2(\frac{\sigma^2}{24}+\frac{|\nu|}6)$. In time,
$\frac1k(u^{n+1}-u^n)=\partial_\tau u^n+\frac k2\partial_{\tau\tau}u(\zeta)$ and
$\theta\partial_\tau u^{n+1}+(1-\theta)\partial_\tau u^n=\partial_\tau u^n+\theta k\partial_{\tau\tau}u(\zeta')$, which differ by at most
$(\frac12+\theta)kM$. For $\theta=\frac12$, expanding about $\tau_n+\frac k2$ instead, the two quantities differ from
$\partial_\tau u(\tau_n+\frac k2)$ by $\frac{k^2}{24}\partial^3_\tau u$ and $\frac{k^2}8\partial^3_\tau u$ at intermediate points.

(iii) The exact values satisfy $Bu^{n+1}=Cu^n+kT^n$, the scheme $BU^{n+1}=CU^n$, so $e^n=u^n-U^n$
satisfies $e^{n+1}=Ge^n+kB^{-1}T^n$ and $e^0=0$. Hence $e^n=k\sum_{m<n}G^{n-1-m}B^{-1}T^m$ and
$\|e^n\|\le nk\,C_s\max_m\|T^m\|$. For the converse see \textcite{strikwerda-finite-difference},
Theorem~1.5.1, proved in Chapter~10.

(iv) $B$ has diagonal $b_0=1+\theta k(\sigma^2/h^2+r)$ and off-diagonal entries $-\theta kp_\pm\le0$ (as
$h|\nu|\le\sigma^2$ means $p_\pm\ge0$). Write $B=b_0(I-N)$, where $N\ge0$ entrywise and
$\|N\|_\infty=\theta k\sigma^2h^{-2}/b_0<1$; the Neumann series $B^{-1}=b_0^{-1}\sum_mN^m$ is positive, with
$\|B^{-1}\|_\infty\le1/(b_0(1-\|N\|_\infty))=1/(1+\theta kr)\le1$. $C$ has entries $(1-\theta)kp_\pm\ge0$ and
$1-(1-\theta)k(\sigma^2/h^2+r)\ge0$, with row sums $1-(1-\theta)kr\le1$. So $G=B^{-1}C$ is positive with
$\|G\|_\infty\le1$. For Crank--Nicolson and $\nu=r=0$, $B+C=2I$, so $G=2B^{-1}-I$. The entries of
$B^{-1}$ off the diagonal are positive (Neumann series) and its diagonal entry is the mean of
$1/(1-\frac12z)=1/(1+2\lambda\sin^2\frac\varphi2)$ over $\varphi=\xi h\in[-\pi,\pi]$:
\[
  \frac1{2\pi}\int_{-\pi}^{\pi}\frac{\dd\varphi}{1+\lambda-\lambda\cos\varphi}=\frac1{\sqrt{(1+\lambda)^2-\lambda^2}}=\frac1{\sqrt{1+2\lambda}},
\]
by the standard integral $\int_{-\pi}^\pi(\alpha-\beta\cos\varphi)^{-1}\dd\varphi=2\pi(\alpha^2-\beta^2)^{-1/2}$, $\alpha>|\beta|$. So the
diagonal entry of $G$ is $2/\sqrt{1+2\lambda}-1$, nonnegative if and only if $\lambda\le\frac32$.
\end{proof}

With (ii), part (iii) gives the rates in the maximum norm whenever the scheme is positive:
$|u^n_j-U^n_j|\le TM[(\frac12+\theta)k+(\frac{\sigma^2}{24}+\frac{|\nu|}6)h^2]$. In $\|\cdot\|_{2,h}$, $\|T^m\|_{2,h}$ is of the same
order when the derivatives of (ii) are also square-integrable in $x$ (write the Taylor remainders
in integral form and use the Cauchy--Schwarz inequality). Von Neumann analysis is in Section~2.2
of \textcite{strikwerda-finite-difference}.

\begin{remark}[Stable is not positive]\label{rem:fd-positivity}
The implicit scheme is positive for every $k$; Crank--Nicolson only for $\lambda$ of order one,
$k\lesssim h^2/\sigma^2$, the very restriction it was chosen to escape. No scheme does better:
\textcite{bolley-crouzeix-1978}, Theorem~1, show that no linear one-step time discretization of
order two or more preserves positivity for all step sizes (their sufficient condition for
Crank--Nicolson is $\lambda\le1$). On a finite grid with zero boundary values the row next to the
boundary binds: there one needs $\lambda\le4-2\sqrt2\approx1.17$ (\exref{ex:fd-theta-positivity}; for $J\ge11$
our grids agree to ten digits). A negative weight is a negative state price. The discrete gamma of a
call at maturity is a spike at the strike, and one Crank--Nicolson step with $\lambda>\frac32$ makes its
centre negative (\cref{fig:fd-schemes}(b)). The other condition, $h|\nu|\le\sigma^2$, fails as $\sigma\to0$, when
the equation becomes a transport equation: the positive alternative to central differences,
upwinding, is only first order, and by Godunov's theorem no linear positive scheme for
transport is better \citep{godunov-1959}. CRR trees meet the same wall at $n\le(r-y)^2T/\sigma^2$
(\exref{ex:fd-upwind}).
\end{remark}

\begin{remark}[Kinks and Rannacher's start]\label{rem:fd-rannacher}
The theorem assumes a smooth solution; a call is not smooth at $(0,\log K)$. Its kink feeds the
highest frequencies, which the implicit scheme damps at each step, but Crank--Nicolson multiplies
the mode $\xi h=\pi$ by $(1-2\lambda)/(1+2\lambda)\to-1$ (with $r=\nu=0$): with $k$ proportional to $h$, the natural choice for a
second-order scheme, $\lambda$ grows like $1/h$ and the sawtooth survives. On the test case with
$k=h/\sigma$, the Crank--Nicolson price error is of first order ($3.8\times10^{-3}$, $2.0\times10^{-3}$, $1.0\times10^{-3}$ for
$J=800,1600,3200$ intervals over $\pm6\sigma\sqrt T$) and the error of its gamma at $S_0$ grows ($-0.15$, $-0.31$,
$-0.60$, while the gamma is $0.0195$). \textcite{rannacher-1984}'s remedy replaces the first Crank--Nicolson
steps by implicit ones. \textcite{giles-carter-2006} prove that four implicit half-steps in
place of the first two steps restore second order in the maximum norm for the value, delta
and gamma of calls and puts; on the grids above the price errors become $-4.0\times10^{-4}$,
$-1.0\times10^{-4}$, $-2.5\times10^{-5}$ and the gamma errors $1.5\times10^{-6}$, $3.9\times10^{-7}$, $9.6\times10^{-8}$. Short
maturities are the worst case: for $T=0.05$ ($J=800$ over $\pm10\sigma\sqrt T$, $N=20$) the plain gamma at $S_0$ is
$2.93$, the Rannacher gamma $0.0892$, the true one $0.0891$.
Placing the strike on a node or midway between two, and averaging discontinuous payoffs over
a cell, complete the treatment \citep{pooley-vetzal-forsyth-2003}; the solvers of
\cref{ch:04-static-replication} use a variant (two implicit steps at each end), and the BBS tree is
its analogue.
\end{remark}

\begin{figure}[ht]
  \centering
  \includegraphics[width=\textwidth]{ch20-schemes}
  \caption{(a) Absolute error of the test-case call at $S_0$ on grids of $J$ intervals over
    $\log S_0\pm6\sigma\sqrt T$: explicit with $\lambda=0.4$, implicit, Crank--Nicolson and Crank--Nicolson with
    Rannacher's start, the last three with $k\approx h/\sigma$. (b) One Crank--Nicolson step of
    $\partial_\tau u=\frac12\sigma^2\partial_{xx}u$ applied to a unit spike at node $0$, for $\lambda=1$, $\frac32$, $4$: the centre
    weight is $2/\sqrt{1+2\lambda}-1$. (c) The gamma at $t=0$ of the test-case call against the spot, on
    the grid $J=800$, $N=25$ ($\lambda=89$): plain Crank--Nicolson (between $-1.49$ and $0.68$ near
    the strike, clipped), with Rannacher's start (on top of the Black--Scholes gamma, within
    $8\times10^{-6}$).}
  \label{fig:fd-schemes}
\end{figure}

\section{Boundaries and the choice of grid}\label{sec:fd-boundary}

A grid must stop somewhere. On $D=(x_\ell,x_u)$ the scheme needs values at the two ends, and the
question is how much a wrong value there costs at the spot. Let $X_s=x+\nu s+\sigma W_s$,
$\vartheta=\inf\{s\ge0:X_s\notin D\}$, and, for bounded measurable boundary data $\beta(\tau,x_\ell)$, $\beta(\tau,x_u)$,
\begin{equation}\label{eq:fd-dirichlet}
  u_D(\tau,x)=\EQ\Big[e^{-r\tau}g(e^{X_\tau})\1_{\{\vartheta>\tau\}}+e^{-r\vartheta}\beta(\tau-\vartheta,X_\vartheta)\1_{\{\vartheta\le\tau\}}\Big],
\end{equation}
the price of the option that pays $\beta$ when the log price leaves $D$ first. When the Dirichlet
problem for \eqref{eq:fd-pde} on $D$ with data $g$ and $\beta$ has a bounded classical solution, Itô's
formula and optional stopping at $\vartheta\wedge\tau$, as in the proof of \cref{thm:sm-feynman-kac}, show that it is
$u_D$; this is the function that the grids of \cref{sec:fd-schemes} approximate.

\begin{proposition}[Far-field boundary conditions]\label{prop:fd-boundary}
Let $g$ be measurable with $|g(s)|\le a(1+s)$, $r\ge0$, $x\in D$, and suppose
$|\beta(t,x_b)-u(t,x_b)|\le\varepsilon$ for $t\in[0,\tau]$ and $x_b\in\{x_\ell,x_u\}$. Then
\[
  |u_D(\tau,x)-u(\tau,x)|\le\varepsilon\,\Q(\vartheta\le\tau)
  \le\varepsilon\Big[e^{-(x_u-x-\nu^+\tau)^2/(2\sigma^2\tau)}+e^{-(x-x_\ell-\nu^-\tau)^2/(2\sigma^2\tau)}\Big],
\]
where the bound holds when $x_u-x\ge\nu^+\tau$ and $x-x_\ell\ge\nu^-\tau$ (a term with $x_\ell=-\infty$ or
$x_u=\infty$ is absent). For the call with the asymptotic values $\beta(t,x_\ell)=0$ and
$\beta(t,x_u)=e^{x_u-yt}-Ke^{-rt}$, $\varepsilon$ is the larger of the call price at $e^{x_\ell}$ and the put price at
$e^{x_u}$ over $t\le\tau$.
\end{proposition}

\begin{proof}
By the Markov property, $M_s=e^{-rs}u(\tau-s,X_s)=\EQ[e^{-r\tau}g(e^{X_\tau})\mid\F_s]$, $s\le\tau$, is a
martingale, closed by an integrable variable (the growth bound and $\EQ[e^{X_\tau}]<\infty$); optional
stopping at the bounded time $\vartheta\wedge\tau$ gives
$u(\tau,x)=\EQ\big[e^{-r\tau}g(e^{X_\tau})\1_{\{\vartheta>\tau\}}+e^{-r\vartheta}u(\tau-\vartheta,X_\vartheta)\1_{\{\vartheta\le\tau\}}\big]$. Subtracting from
\eqref{eq:fd-dirichlet}, $u_D-u=\EQ\big[e^{-r\vartheta}(\beta-u)(\tau-\vartheta,X_\vartheta)\1_{\{\vartheta\le\tau\}}\big]$, which is at most
$\varepsilon\,\Q(\vartheta\le\tau)$ in absolute value since $r\ge0$. The path leaves through $x_u$ before $\tau$ only if
$\max_{s\le\tau}\sigma W_s\ge x_u-x-\nu^+\tau=:m$, which by the reflection principle
(\cref{thm:bar-reflection}) has probability $2N(-m/(\sigma\sqrt\tau))\le e^{-m^2/(2\sigma^2\tau)}$; the last
inequality holds for $m\ge0$ because $e^{-z^2/2}-2N(-z)$ vanishes at $0$ and at $\infty$ and increases and
then decreases on $[0,\infty)$. The lower end is symmetric. For the call, $u-\beta$ at $x_u$ is the put
price, by put--call parity.
\end{proof}

The bound is Gaussian in the distance to the boundary in units of $\sigma\sqrt\tau$, and it multiplies
the error of the boundary data: on $\pm m\sigma\sqrt T$ around the spot even the crude choice $\beta=g$
costs at most about $\max|g-u|\cdot2e^{-m^2/2}$, $3\times10^{-8}\max|g-u|$ for $m=6$, the default of our code.
\Cref{fig:fd-american}(b) measures the truncation error of the Crank--Nicolson grid. For positive
schemes the same argument runs on the Markov chain of the scheme, whose steps are kernels with
nonnegative weights summing to at most one; \textcite{kangro-nicolaides-2000} prove the
exponential decay for the Black--Scholes equation itself.

\emph{The course's example.} The PINN of the course's discussion notebook prices the call with
$K=40$, $r=7\%$, $\sigma=60\%$, $T=1$ on $S\in[0,120]$ with the value $S_{\max}-K=80$ at the top, whereas the
asymptotic value is $82.70$ and the price $83.08$. On a Crank--Nicolson grid in $S$ ($1200$ intervals,
$600$ steps) the value $80$ makes the price too low by $0.036$ at $S_0=40$ and by $0.75$ at $S_0=80$; the
asymptotic value reduces these to $0.0007$ and $0.050$. The proposition, with $\varepsilon=3.08$ and $0.378$
and the exact exit probabilities $0.047$ and $0.44$, bounds them by $0.15$, $1.35$, $0.018$ and $0.17$.
A grid in $S$ needs no condition at $S=0$, where the equation reduces to $\partial_\tau V=-rV$, but its
local $\lambda$ grows like $S^2$, so its explicit version needs $k\lesssim h_S^2/(\sigma S_{\max})^2$ for a mesh $h_S$; in log
price $\lambda$ is the same everywhere.

\section{American options on a grid}\label{sec:fd-american}

The American price is the value of an optimal stopping problem (\cref{ch:14-american-mc}), and
it solves the linear complementarity problem
$u\ge g$, $\partial_\tau u-\mathcal Lu\ge0$, $(u-g)(\partial_\tau u-\mathcal Lu)=0$, with $\mathcal L$ the right side of \eqref{eq:fd-pde}
\citep{jaillet-lamberton-lapeyre-1990}. Each step of the $\theta$-scheme becomes a finite-dimensional
complementarity problem: given $q=CV^n+kF^n$ and the payoff vector $g$, find $z=V^{n+1}$ with
\begin{equation}\label{eq:fd-lcp}
  z\ge g,\qquad Bz-q\ge0,\qquad(z-g)\cdot(Bz-q)=0,
\end{equation}
inequalities entry by entry. With $\Phi_j(z)=\max\big(g_j,(q_j-\sum_{l\ne j}B_{jl}z_l)/B_{jj}\big)$,
\emph{projected SOR} (PSOR, \textcite{cryer-1971}) sweeps through the nodes with a relaxation
$\omega\in(0,2)$, $z_j\leftarrow\max\big(g_j,z_j+\omega(\Phi^0_j(z)-z_j)\big)$, where $\Phi^0_j$ is $\Phi_j$ without the maximum and
uses the entries already updated; $\omega=1$ is \emph{projected Gauss--Seidel},
$z_j\leftarrow\Phi_j(z)$. The algorithm of \textcite{brennan-schwartz-1977} eliminates the
super-diagonal of $B$ from the top, $d'_{J-1}=B_{J-1,J-1}$, $q'_{J-1}=q_{J-1}$ and, for
$j=J-2,\dots,1$,
\begin{equation}\label{eq:fd-bs}
  m_j=\frac{B_{j,j+1}}{d'_{j+1}},\quad d'_j=B_{jj}-m_jB_{j+1,j},\quad q'_j=q_j-m_jq'_{j+1},
\end{equation}
and then, upwards for $j=1,\dots,J-1$ with $B_{1,0}z_0:=0$,
$z_j=\max\big(g_j,(q'_j-B_{j,j-1}z_{j-1})/d'_j\big)$: one elimination and one projected substitution,
$O(J)$ operations.

\begin{proposition}[American options on a grid]\label{prop:fd-american}
Let $\theta\in(0,1]$, $r\ge0$, $h|\nu|\le\sigma^2$, and let $B=I-\theta kA$ be the matrix of one step on the
interior nodes $1,\dots,J-1$. Then $B_{jl}\le0$ for $l\ne j$ and every row sum of $B$ is at least
$1+\theta kr$.
\begin{enumerate}[label=(\roman*)]
  \item For all $q,g$, \eqref{eq:fd-lcp} has exactly one solution $z^*$. It is the smallest
    \emph{supersolution}: $z^*\le s$ for every $s$ with $s\ge g$ and $Bs\ge q$; and $z^*=\Phi(z^*)$.
  \item Projected Gauss--Seidel started at $z^{(0)}=g$ produces a nondecreasing sequence, bounded by
    every supersolution, that converges to $z^*$, for every order of the sweep. If $h|\nu|<\sigma^2$,
    PSOR converges to $z^*$ for every $\omega\in(0,2)$ and every starting point.
  \item If $z^*$ has its exercise region at the bottom, that is, for some $j^*$, $z^*_j=g_j$ for
    $j\le j^*$ and $(Bz^*-q)_j=0$ for $j>j^*$, then \eqref{eq:fd-bs} returns $z^*$. The mirror image,
    eliminating from the bottom, holds for an exercise region at the top.
  \item For $\rho>0$, the penalized equation $Bz=q+\rho(g-z)^+$ has a unique solution $z_\rho$, the
    policy iteration of \textcite{forsyth-vetzal-2002} finds it in finitely many steps (their
    Theorem~6.1), and $z_\rho\le z^*\le z_\rho+\|(Bg-q)^+\|_\infty/\rho$.
\end{enumerate}
\end{proposition}

\begin{proof}
The off-diagonal entries are $-\theta kp_\pm\le0$, and the row sums are $1+\theta kr$ for interior rows and
larger for the first and the last, which miss a neighbour.

(i) A solution $z$ satisfies $z_j>g_j\Rightarrow(Bz)_j=q_j$ and $z_j=g_j\Rightarrow(Bz)_j\ge q_j$; dividing by $B_{jj}>0$, this
is $z_j=\Phi_j(z)$, and conversely. Let $s$ be a supersolution and suppose $w=z-s$ has a positive
maximum at $j_0$. Then $z_{j_0}>s_{j_0}\ge g_{j_0}$, so $(Bz)_{j_0}=q_{j_0}\le(Bs)_{j_0}$ and $(Bw)_{j_0}\le0$; but
$(Bw)_{j_0}\ge(\sum_lB_{j_0l})w_{j_0}>0$ because $B_{j_0l}\le0$ and $w_l\le w_{j_0}$. So $z\le s$. A solution is a
supersolution, so two solutions are each below the other: there is at most one, and it is the
smallest supersolution. Existence follows from (ii).

(ii) $\Phi$ is continuous and nondecreasing in each entry, because $B_{jl}\le0$. A sweep replaces the
entries one by one by $\Phi_j$ of the current vector, so the sweep map is nondecreasing too. Since
$\Phi(g)\ge g$, the first sweep does not decrease any entry, and by induction the sequence is
nondecreasing. If $s$ is a supersolution, $\Phi(s)\le s$ entrywise, so $z^{(i)}\le s$ implies
$z^{(i+1)}\le s$; and $s=c\mathbf 1$ is one as soon as $c\ge\max(\max_jg_j,\max_jq_j,0)$, since $B\mathbf 1\ge\mathbf 1$. The
bounded nondecreasing sequence converges to a fixed point of $\Phi$, which solves
\eqref{eq:fd-lcp}. For PSOR: if $p_\pm>0$, $D=\operatorname{diag}((p_-/p_+)^{j/2})$ makes $P=D^{-1}BD$ symmetric, with
off-diagonal entries $-\theta k\sqrt{p_-p_+}$, and positive definite by Gershgorin, since
$2\sqrt{p_-p_+}\le p_-+p_+$. PSOR for $B$, $q$, $g$ is PSOR for $P$, $D^{-1}q$, $D^{-1}g$, which converges for
$\omega\in(0,2)$ by \textcite{cryer-1971}.

(iii) By induction from the top, $d'_j\ge1+|B_{j,j-1}|$ and $m_j\le0$: $d'_{J-1}=B_{J-1,J-1}$, and if
$d'_{j+1}\ge1+|B_{j+1,j}|$, then $0\le m_jB_{j+1,j}\le|B_{j,j+1}|$, so $d'_j\ge B_{jj}-|B_{j,j+1}|\ge1+|B_{j,j-1}|$.
Row $j$ of the eliminated system, $B_{j,j-1}z_{j-1}+d'_jz_j=q'_j$, is row $j$ of $Bz=q$ plus $|m_j|$ times
row $j+1$ of the eliminated system. So the residuals $r'_j=B_{j,j-1}z^*_{j-1}+d'_jz^*_j-q'_j$ satisfy
$r'_{J-1}=r_{J-1}$ and $r'_j=r_j+|m_j|r'_{j+1}$, where $r=Bz^*-q\ge0$ vanishes above $j^*$. Hence $r'_j=0$ for
$j>j^*$ and $r'_j\ge0$ for $j\le j^*$. For $j>j^*$ this says $z^*_j=(q'_j-B_{j,j-1}z^*_{j-1})/d'_j$, which is at least $g_j$;
for $j\le j^*$, $z^*_j=g_j\ge(q'_j-B_{j,j-1}z^*_{j-1})/d'_j$. In both cases
$z^*_j=\max\big(g_j,(q'_j-B_{j,j-1}z^*_{j-1})/d'_j\big)$, so the upward recursion, started at $j=1$, reproduces $z^*$.

(iv) Existence, uniqueness and termination are Theorem~6.1 of \textcite{forsyth-vetzal-2002}; the
bracket is \exref{ex:fd-penalty}.
\end{proof}

The hypothesis of (iii) is not automatic. For the put in the implicit scheme it is justified
by \textcite{jaillet-lamberton-lapeyre-1990}; our code checks the residual of \eqref{eq:fd-lcp}: below
$10^{-12}$ for the options of \cref{tab:fd-american}, but $5\times10^{-3}$ and $7\times10^{-3}$, with the exercise region
assumed at the top and at the bottom, for an American straddle with $r=y=5\%$, which has an exercise region at
each end (\exref{ex:fd-straddle}). Bermudan exercise
replaces the complementarity problem by a maximum after the linear solve at the dates. Convergence of the discrete American prices to the American price holds
for positive schemes by the theory of monotone schemes of \textcite{barles-souganidis-1991}, and
for the implicit scheme by \textcite{jaillet-lamberton-lapeyre-1990}; Crank--Nicolson on our grids,
with $\lambda$ from $3$ to $23$, is outside both, and its convergence is observed, not proved here.

\begin{table}[ht]
\centering
\footnotesize
\setlength{\tabcolsep}{4pt}
\begin{tabular}{@{}lccccccc@{}}
\toprule
& European & grid & corrected & BBSR & LR & trinomial & Ch.~14\\
\midrule
test case, put & $8.349406$ & $8.349393$ & $8.349408$ & $8.349408$ & $8.349408$ & $8.349154$ & $8.349408$\\
test case, call & $7.364290$ & $7.459694$ & $7.459709$ & $7.459739$ & $7.459766$ & $7.459608$ & $7.4597$\\
LS $36$, $0.2$, $1$ & $3.844308$ & $4.486662$ & $4.486666$ & $4.486679$ & $4.486560$ & $4.486642$ & $4.487$\\
LS $40$, $0.2$, $1$ & $2.066401$ & $2.319556$ & $2.319562$ & $2.319579$ & $2.319535$ & $2.319522$ & $2.320$\\
LS $44$, $0.4$, $2$ & $5.201995$ & $5.646682$ & $5.646696$ & $5.646747$ & $5.646698$ & $5.646838$ & $5.647$\\
\bottomrule
\end{tabular}
\caption{American puts and calls. Grid: Crank--Nicolson with Rannacher's start, $J=3200$
  intervals over $\pm6\sigma\sqrt T$, $N=1600$ steps, Brennan--Schwartz (PSOR and the penalty method
  agree to $2\times10^{-6}$); corrected: plus the Black--Scholes price minus the European price on the
  same grid. BBSR: $n=4000$ (the tree of \cref{ch:14-american-mc}, which used $2000$ steps for
  \cref{tab:am-ls}); LR: Leisen--Reimer,
  $n=4001$; trinomial: $\kappa=\sqrt{3/2}$, $n=2000$. LS $S_0$, $\sigma$, $T$: the puts of Longstaff and
  Schwartz, $K=40$, $r=6\%$, $y=0$. \Cref{ch:14-american-mc} quotes the test-case values and
  \cref{tab:am-ls} to the digits shown.}
\label{tab:fd-american}
\end{table}

\paragraph{The numbers.} \Cref{tab:fd-american} reproduces the converged American values of
\cref{ch:14-american-mc}: the test-case put $8.349408$ (early-exercise premium $2.0\times10^{-6}$ on the
grid, as in \cref{ch:14-american-mc}), the test-case call $7.4597$ ($y>r$, so early exercise of the call pays),
and the Longstaff--Schwartz rows. With exercise at $50$ dates per year the grid gives the
Bermudan values of \cref{ch:14-american-mc} within $10^{-4}$ ($4.477807$ against $4.477810$ for the first LS row).
Brennan--Schwartz costs one pass, the penalty method about $1.1$ tridiagonal solves per step on
average (at most $5$), and PSOR with $\omega=1.2$ up to $16$, $36$ and $70$ sweeps per step ($15$, $34$, $65$ on
average) for $J=400$, $800$, $1600$ with $N=J/2$: an iterative solver slows down as $\lambda$ grows. The American errors fall by
about $3$ when $J$ and $N$ double (order about $1.6$), against $4$ for the European ones: the price is
not twice differentiable across the exercise boundary of \cref{fig:fd-american}(a).

\begin{figure}[ht]
  \centering
  \includegraphics[width=\textwidth]{ch20-american}
  \caption{(a) Exercise boundary of the Longstaff--Schwartz put ($S_0=36$, $K=40$, $r=6\%$, $\sigma=20\%$,
    $T=1$): grid ($J=1600$, $N=800$, Brennan--Schwartz; largest exercised node) and CRR tree ($2000$
    steps, every tenth shown); $K$ dashed, $L^*=2rK/(2r+\sigma^2)=30$ dotted. (b) Test-case call: the
    Crank--Nicolson price ($k=1/400$, $h=0.003$) on $\log S_0\pm m\sigma\sqrt T$ minus that on $\log S_0\pm10\sigma\sqrt T$,
    for asymptotic and payoff boundary values, and the bounds of \cref{prop:fd-boundary} (dashed);
    differences below $10^{-11}$ (rounding) are not shown.}
  \label{fig:fd-american}
\end{figure}

\section{Path dependence and several factors}\label{sec:fd-extensions}

\paragraph{Barriers.} A continuously monitored knock-out is a Dirichlet condition: the grid ends
at the barrier, a node, with value zero. A discretely monitored one sets the values beyond the
barrier to zero at the monitoring dates; we place the barrier midway between two nodes and
repeat Rannacher's start after each date. On the course's monthly down-and-out call ($B=95$)
the Crank--Nicolson grid with $J=1600$, $N=960$ gives $5.32995$ against the exact $5.32979$ of \cref{ch:16-barriers}, and $3.75411$ against $3.75412$ for
continuous monitoring (compare \cref{prop:bar-bgk}). In a tree the error oscillates in $n$ unless
a layer of nodes lies on the barrier \citep{boyle-lau-1994}.

\paragraph{Other path dependence and several factors.} A running maximum or average becomes an
extra state variable. In two dimensions (Heston,
\cref{prop:sv-pde}; two assets) the implicit step is a sparse system rather than a tridiagonal one;
alternating-direction implicit methods split it into tridiagonal solves along each axis
\citep{peaceman-rachford-1955}, with the mixed derivative of correlated factors treated
explicitly \citep{craig-sneyd-1988,in-t-hout-foulon-2010}. For a target error $\varepsilon$ and a smooth
solution, the explicit scheme needs $h\approx\varepsilon^{1/2}$ and $k\approx h^2$, Crank--Nicolson with Rannacher's
start $h\approx k\approx\varepsilon^{1/2}$; in $d$ space dimensions they cost $\varepsilon^{-(d+2)/2}$ and $\varepsilon^{-(d+1)/2}$
operations, the course's rule $\varepsilon^{-(d+1)/\alpha}$ for lattices with $\alpha=2$, against $\varepsilon^{-2}$ for Monte
Carlo (\cref{prop:mc-dimension}).

\begin{incode}
\coderef{code/src/mfe230d/grids.py} implements the trees, \texttt{fd\_solve} (the $\theta$-scheme with
barriers and American exercise by \texttt{psor}, \texttt{brennan\_schwartz} or \texttt{penalty}) and the
stability tools; \coderef{code/tests/test_grids.py} checks every number quoted here.
\end{incode}

\demo{20-grids.ipynb}{trees, stability, the far field and the American put.}

\furtherreading
Finance: \textcite{wilmott-howison-dewynne}; stability: \textcite{strikwerda-finite-difference}.
